% ============================================================
\chapter{Martingales --- Temps Continu}
\label{ch:martingales_continu}
% ============================================================

Le passage des martingales en temps discret au temps continu est l'une des transitions les plus d\'elicates de la th\'eorie des processus stochastiques. En temps discret, on conditionne par rapport \`a un pass\'e fini~; en temps continu, la filtration $(\mathcal{F}_t)_{t \geq 0}$ forme un continuum d'informations, et des ph\'enom\`enes subtils apparaissent~: r\'egularit\'e des trajectoires, temps d'arr\^et non d\'enombrables, in\'egalit\'es maximales. Joseph Doob, dans les ann\'ees 1950, a montr\'e que toute martingale en temps continu admet une version \`a trajectoires c\`adl\`ag (continues \`a droite, avec limites \`a gauche), un r\'esultat technique mais essentiel. Les in\'egalit\'es de Doob, le th\'eor\`eme d'arr\^et optionnel, et la d\'ecomposition de Doob--Meyer forment le socle sur lequel repose tout le calcul stochastique --- et, par extension, la finance math\'ematique moderne.

\section{D\'efinitions}

\begin{definition}[Martingale en temps continu]
Soit $(\Omega, \mathcal{F}, (\mathcal{F}_t)_{t \geq 0}, \PP)$ un espace de
probabilit\'e filtr\'e satisfaisant les conditions habituelles. Un processus
adapt\'e $(M_t)_{t \geq 0}$ est une \emph{martingale} (en temps continu) si~:
\begin{enumerate}
  \item $\E[\abs{M_t}] < \infty$ pour tout $t \geq 0$,
  \item $\E[M_t \mid \mathcal{F}_s] = M_s$ p.s.\ pour tous $0 \leq s \leq t$.
\end{enumerate}
\end{definition}

\begin{definition}[Martingale locale]
Un processus adapt\'e continu $(M_t)_{t \geq 0}$ est une \emph{martingale locale}
s'il existe une suite croissante de temps d'arr\^et $(\tau_n)_{n \geq 1}$ avec
$\tau_n \uparrow \infty$ p.s.\ telle que pour chaque $n$, le processus arr\^et\'e
$(M_{t \wedge \tau_n})_{t \geq 0}$ est une martingale.

La suite $(\tau_n)$ est appel\'ee \emph{suite localisante}.
\end{definition}

\begin{attention}[title=\textbf{Martingale locale $\neq$ martingale}]
Toute martingale est une martingale locale, mais la r\'eciproque est fausse.
Une martingale locale positive est une sur-martingale (lemme de Fatou), mais pas
n\'ecessairement une martingale.
\end{attention}

\section{R\'egularit\'e des trajectoires}

\begin{theoreme}[Modification c\`adl\`ag]
\label{thm:cadlag}
Soit $(M_t)_{t \geq 0}$ une martingale (ou sur-martingale) \`a valeurs r\'eelles.
Si $t \mapsto \E[M_t]$ est continue \`a droite, alors $(M_t)$ admet une
modification c\`adl\`ag.
\end{theoreme}

\begin{remarque}
Sous les conditions habituelles de la filtration, la modification c\`adl\`ag est
encore une martingale. Dor\'enavant, nous supposerons que les martingales en
temps continu sont c\`adl\`ag.
\end{remarque}

\begin{theoreme}[Continuit\'e des martingales $L^2$]
Si $(M_t)$ est une martingale de carr\'e int\'egrable dont les trajectoires sont
continues, alors le processus de variation quadratique $\langle M \rangle_t$ est
continu et $(M_t^2 - \langle M \rangle_t)$ est une martingale.
\end{theoreme}

\section{In\'egalit\'es de Doob en temps continu}

\begin{theoreme}[In\'egalit\'e maximale de Doob]
Soit $(M_t)_{0 \leq t \leq T}$ une martingale c\`adl\`ag. Pour tout $\lambda > 0$~:
\[
  \PP\Bigl(\sup_{0 \leq t \leq T} \abs{M_t} \geq \lambda\Bigr)
  \leq \frac{\E[\abs{M_T}]}{\lambda}.
\]
\end{theoreme}

\begin{theoreme}[In\'egalit\'e $L^p$ de Doob en temps continu]
Pour $p > 1$ et une martingale c\`adl\`ag $(M_t)$~:
\[
  \E\Bigl[\sup_{0 \leq t \leq T} \abs{M_t}^p\Bigr]
  \leq \Bigl(\frac{p}{p-1}\Bigr)^p \E[\abs{M_T}^p].
\]
\end{theoreme}

\begin{theoreme}[In\'egalit\'e de Burkholder--Davis--Gundy]
\label{thm:BDG}
Il existe des constantes universelles $c_p, C_p > 0$ telles que pour toute
martingale locale continue $(M_t)$ avec $M_0 = 0$ et pour tout $p > 0$~:
\[
  c_p \, \E[\langle M \rangle_T^{p/2}]
  \leq \E\Bigl[\sup_{0 \leq t \leq T} \abs{M_t}^p\Bigr]
  \leq C_p \, \E[\langle M \rangle_T^{p/2}].
\]
Pour $p = 2$, on peut prendre $c_2 = 1$ et $C_2 = 4$ (in\'egalit\'e de Doob).
\end{theoreme}

\section{Variation quadratique}

\begin{definition}[Variation quadratique]
\label{def:variation_quadratique}
Soit $(M_t)_{t \geq 0}$ une martingale locale continue. La \emph{variation
quadratique} (ou \emph{crochet}) de $M$ est l'unique processus continu croissant
$(\langle M \rangle_t)_{t \geq 0}$ avec $\langle M \rangle_0 = 0$ tel que
$(M_t^2 - \langle M \rangle_t)$ soit une martingale locale.

On a la repr\'esentation~:
\[
  \langle M \rangle_t = \lim_{\abs{\Pi} \to 0}
  \sum_{k=0}^{n-1} (M_{t_{k+1}} - M_{t_k})^2 \quad \text{(limite en probabilit\'e)},
\]
o\`u $\Pi = \{0 = t_0 < t_1 < \cdots < t_n = t\}$ est une partition de $[0, t]$.
\end{definition}

\begin{definition}[Covariation]
Pour deux martingales locales continues $M$ et $N$, la \emph{covariation}
(ou \emph{crochet crois\'e}) est
\[
  \langle M, N \rangle_t = \frac{1}{4}\bigl(\langle M + N \rangle_t
  - \langle M - N \rangle_t\bigr),
\]
et c'est l'unique processus continu \`a variation finie tel que
$(M_t N_t - \langle M, N \rangle_t)$ soit une martingale locale.
\end{definition}

\begin{formules}[title=\textbf{Variations quadratiques fondamentales}]
Pour le mouvement brownien $(B_t)$~:
\begin{align*}
  \langle B \rangle_t &= t \\
  \langle B^{(i)}, B^{(j)} \rangle_t &= \delta_{ij} t
    \quad \text{(brownien multidimensionnel)}
\end{align*}
\end{formules}

\section{D\'ecomposition de Doob--Meyer}

\begin{theoreme}[D\'ecomposition de Doob--Meyer]
\label{thm:doob_meyer}
Soit $(X_t)_{t \geq 0}$ une sur-martingale c\`adl\`ag de classe (D). Alors il
existe une d\'ecomposition unique~:
\[
  X_t = M_t - A_t,
\]
o\`u $(M_t)$ est une martingale c\`adl\`ag et $(A_t)$ est un processus croissant
pr\'evisible avec $A_0 = 0$.
\end{theoreme}

\begin{definition}[Classe (D)]
Une sur-martingale $(X_t)$ est de \emph{classe (D)} si la famille
$\{X_{\tau} : \tau \text{ temps d'arr\^et fini}\}$ est uniform\'ement int\'egrable.
\end{definition}

\begin{remarque}
La d\'ecomposition de Doob--Meyer est l'analogue en temps continu de la
d\'ecomposition de Doob en temps discret. Elle est fondamentale pour la
construction de l'int\'egrale stochastique.
\end{remarque}

\section{Th\'eor\`emes d'arr\^et en temps continu}

\begin{theoreme}[Th\'eor\`eme d'arr\^et optionnel --- temps continu]
\label{thm:arret_continu}
Soit $(M_t)_{t \geq 0}$ une martingale continue uniform\'ement int\'egrable et
$\tau$ un temps d'arr\^et (\'eventuellement infini). Alors
\[
  \E[M_{\tau}] = \E[M_0].
\]
Plus g\'en\'eralement, si $\sigma \leq \tau$ sont deux temps d'arr\^et~:
\[
  \E[M_{\tau} \mid \mathcal{F}_{\sigma}] = M_{\sigma} \quad \text{p.s.}
\]
\end{theoreme}

\begin{theoreme}[Th\'eor\`eme d'arr\^et pour les sur-martingales positives]
\label{thm:arret_surmartingale}
Si $(M_t)_{t \geq 0}$ est une sur-martingale positive continue et $\tau$ est
un temps d'arr\^et, alors
\[
  \E[M_{\tau}] \leq \E[M_0].
\]
\end{theoreme}

\section{Repr\'esentation des martingales}

\begin{theoreme}[Th\'eor\`eme de repr\'esentation des martingales browniennes]
\label{thm:representation}
Soit $(B_t)_{t \geq 0}$ un mouvement brownien et $(\mathcal{F}_t)$ sa filtration
naturelle compl\'et\'ee. Toute martingale locale continue $(M_t)$ par rapport \`a
$(\mathcal{F}_t)$ admet la repr\'esentation~:
\[
  M_t = M_0 + \int_0^t H_s \, dB_s,
\]
o\`u $(H_t)$ est un processus pr\'evisible avec $\int_0^t H_s^2 \, ds < \infty$
p.s.\ pour tout $t$.
\end{theoreme}

\begin{intuition}[title=\textbf{Compl\'etude de la filtration brownienne}]
Ce th\'eor\`eme exprime que la filtration brownienne est <<~compl\`ete~>> au sens
o\`u toute martingale peut s'\'ecrire comme int\'egrale stochastique par rapport
au brownien. C'est un r\'esultat central pour les applications en finance
(th\'eor\`eme de compl\'etude des march\'es).
\end{intuition}

\section{Th\'eor\`eme de Girsanov}

\begin{theoreme}[Girsanov]
\label{thm:girsanov}
Soit $(B_t)_{t \geq 0}$ un mouvement brownien sous $\PP$ et $(\theta_t)_{t \geq 0}$
un processus adapt\'e avec $\int_0^T \theta_s^2 \, ds < \infty$ p.s. D\'efinissons
l'exponentielle de Dol\'eans--Dade~:
\[
  Z_t = \exp\Bigl(-\int_0^t \theta_s \, dB_s - \frac{1}{2} \int_0^t \theta_s^2 \, ds\Bigr).
\]
Si $(Z_t)$ est une martingale (condition de Novikov~:
$\E[\exp(\frac{1}{2}\int_0^T \theta_s^2 \, ds)] < \infty$), alors sous la mesure
$\mathbb{Q}$ d\'efinie par $d\mathbb{Q}/d\PP\big|_{\mathcal{F}_T} = Z_T$, le
processus
\[
  \tilde{B}_t = B_t + \int_0^t \theta_s \, ds
\]
est un mouvement brownien.
\end{theoreme}

\section{Convergence en temps continu}

\begin{theoreme}[Convergence des martingales continues]
\label{thm:convergence_continue}
Soit $(M_t)_{t \geq 0}$ une sur-martingale continue \`a droite avec
$\sup_{t \geq 0} \E[M_t^+] < \infty$. Alors $M_{\infty} = \lim_{t \to \infty} M_t$
existe p.s.\ et $M_{\infty} \in L^1$.
\end{theoreme}

\section{Exercices}

\begin{exercice}
Soit $(B_t)_{t \geq 0}$ un mouvement brownien. V\'erifier que
$M_t = \exp(\theta B_t - \theta^2 t/2)$ est une martingale pour tout
$\theta \in \R$.
\end{exercice}

\begin{exercice}
D\'emontrer que si $(M_t)$ est une martingale locale continue positive, alors
c'est une sur-martingale. \emph{Indication~: utiliser le lemme de Fatou.}
\end{exercice}

\begin{exercice}
Soit $(M_t)$ une martingale continue de carr\'e int\'egrable. Montrer que
$\E[M_t^2] = \E[M_0^2] + \E[\langle M \rangle_t]$.
\end{exercice}

\begin{exercice}
Calculer $\langle M \rangle_t$ pour $M_t = B_t^2 - t$ o\`u $(B_t)$ est un
mouvement brownien.
\end{exercice}

\begin{exercice}
Soit $(B_t)$ un mouvement brownien et $\tau_a = \inf\{t \geq 0 : B_t = a\}$.
Utiliser le th\'eor\`eme d'arr\^et avec la martingale exponentielle pour calculer
$\E[e^{-\alpha \tau_a}]$ pour $\alpha > 0$ et $a > 0$.
\end{exercice}

\begin{exercice}
\'Enoncer et d\'emontrer le th\'eor\`eme de Girsanov dans le cas o\`u $\theta_t = \theta$
est constant. V\'erifier la condition de Novikov.
\end{exercice}

\begin{exercice}
Montrer que l'in\'egalit\'e de Burkholder--Davis--Gundy (th\'eor\`eme~\ref{thm:BDG})
implique l'in\'egalit\'e $L^2$ de Doob pour les martingales continues.
\end{exercice}
